\begin{definition}[Prime cyclic extension package]\label{mod:ramification-primecyclicextension}
Let $K/F$ be a finite valuative extension of nonarchimedean local fields.
The theorem-level predicate \texttt{CyclicPrimeExtension} already records
exactly the three field-theoretic hypotheses used in the cyclic-prime
argument:
\[
 K/F\text{ is Galois},\qquad
 \operatorname{Gal}(K/F)\text{ is cyclic},\qquad
 \ell=[K:F]\text{ is prime}.
\]
The structure \texttt{PrimeCyclicExtension} reuses that predicate as its
sole field, with constructors and an equivalence in both directions; it
does not unpack and duplicate the extension data.  Here the degree is
definitionally the canonical value
$\ell=\operatorname{finrank}_F K$.  The local fields, algebra structure,
valuative extension, and finite-module structure remain ambient typeclass
data.  In particular, the package contains no independently chosen degree,
Galois generator, ramification invariant, lower filtration, or break.

The Galois and cyclic hypotheses are exposed as instances.  Finite Galois theory
then gives
\[
 \#\operatorname{Gal}(K/F)=\ell,
\]
so the Galois group has prime cardinality.  Cyclicity supplies a
noncomputably chosen automorphism $\sigma_0$ for which
\[
 \langle\sigma_0\rangle=\operatorname{Gal}(K/F),\qquad
 \forall\tau\ \exists n\in\mathbf Z,\quad \sigma_0^n=\tau,
 \qquad \operatorname{ord}(\sigma_0)=\ell.
\]
Consequently $\sigma_0\ne1$, the top and bottom subgroups are distinct,
and every subgroup of the Galois group is either trivial or the whole group.

No extension invariants are duplicated.  The package uses the canonical
quantities
\[
 e=e(K/F)=\texttt{ramificationIndex}(F,K),\qquad
 f=f(K/F)=\texttt{residueDegree}(F,K).
\]
They are positive and satisfy the already established local-field identity
\[
 \ell=ef.
\]
This node draws no prime-factor dichotomy from that identity.  The residue
characteristic is likewise the canonical abbreviation
\[
 p_F=\operatorname{char}k_F
     =\texttt{ringChar}(\operatorname{ResidueField}(F)).
\]
It is prime, and the injective residue-field algebra map proves
$p_K=p_F$ for every valuative extension.

The lower filtration is always the canonical
$G_i=\texttt{lowerRamificationGroup}(F,K,i)$ from the preceding node.
Prime cardinality implies, for every $i\in\mathbf Z$,
\[
 G_i=1\quad\text{or}\quad G_i=G.
\]
For a manuscript-style nonnegative index $t$, define only the minimal
boundary predicate
\[
 \operatorname{IsLowerBreak}(t)
 \quad\Longleftrightarrow\quad
 G_t=G\ \text{ and }\ G_{t+1}=1.
\]
The complete one-break profile is derived rather than stored:
\[
 G_i=G\Longleftrightarrow i\le t,
 \qquad
 G_i=1\Longleftrightarrow t<i.
\]
In particular such a $t$ is unique, and the chosen generator lies in the
exact shell $G_t\setminus G_{t+1}$.

The existence theorem deliberately has the explicit premise $G_0=G$.
Using eventual triviality from the lower-group node, choose the least
$n\in\mathbf N$ with $G_n=1$ and put $t=n-1$.  The premise gives $n>0$,
and the prime-order subgroup dichotomy gives $G_t=G$; hence there is a
unique lower break.  Lean also exposes the resulting canonical selected
index and proves its boundary property and uniqueness.  The implication
from ordinary ramifiedness (or total ramification) to $G_0=G$ is not assumed
here.

Thus this node supplies the raw context and elementary one-break interfaces
needed after \ref{thm:first-main} and \ref{eq:one-break-new}, while leaving
the following results to their designated downstream nodes: the
unramified--totally-ramified and tame--wild dichotomies, the wild equality
$\ell=p_F$, the Herbrand functions and their piecewise formulas, the
different and trace-ideal formulas, and every graded or full norm-filtration
statement from \ref{thm:graded-norm} and \ref{thm:norm-filtration}.
\LeanFile{LanglandsFirstMainLemma/Ramification/PrimeCyclicExtension.lean}
\lean{LanglandsFirstMainLemma.PrimeCyclicExtension}
\lean{LanglandsFirstMainLemma.PrimeCyclicExtension.toCyclicPrimeExtension}
\lean{LanglandsFirstMainLemma.PrimeCyclicExtension.ofCyclicPrimeExtension}
\lean{LanglandsFirstMainLemma.PrimeCyclicExtension.iff_cyclicPrimeExtension}
\lean{LanglandsFirstMainLemma.PrimeCyclicExtension.instIsGalois}
\lean{LanglandsFirstMainLemma.PrimeCyclicExtension.instIsCyclic}
\lean{LanglandsFirstMainLemma.PrimeCyclicExtension.degree_prime}
\lean{LanglandsFirstMainLemma.PrimeCyclicExtension.galoisCard_eq_degree}
\lean{LanglandsFirstMainLemma.PrimeCyclicExtension.galoisCard_prime}
\lean{LanglandsFirstMainLemma.PrimeCyclicExtension.instGaloisCardPrimeFact}
\lean{LanglandsFirstMainLemma.PrimeCyclicExtension.generator}
\lean{LanglandsFirstMainLemma.PrimeCyclicExtension.generator_mem_zpowers}
\lean{LanglandsFirstMainLemma.PrimeCyclicExtension.zpowers_generator}
\lean{LanglandsFirstMainLemma.PrimeCyclicExtension.exists_generator_zpow}
\lean{LanglandsFirstMainLemma.PrimeCyclicExtension.orderOf_generator}
\lean{LanglandsFirstMainLemma.PrimeCyclicExtension.generator_ne_one}
\lean{LanglandsFirstMainLemma.PrimeCyclicExtension.galoisSubgroup_top_ne_bot}
\lean{LanglandsFirstMainLemma.PrimeCyclicExtension.subgroup_eq_bot_or_eq_top}
\lean{LanglandsFirstMainLemma.PrimeCyclicExtension.ramificationIndex_positive}
\lean{LanglandsFirstMainLemma.PrimeCyclicExtension.residueDegree_positive}
\lean{LanglandsFirstMainLemma.PrimeCyclicExtension.degree_eq_ramificationIndex_mul_residueDegree}
\lean{LanglandsFirstMainLemma.residueCharacteristic}
\lean{LanglandsFirstMainLemma.residueCharacteristic_prime}
\lean{LanglandsFirstMainLemma.residueCharacteristic_extension_eq}
\lean{LanglandsFirstMainLemma.PrimeCyclicExtension.IsLowerBreak}
\lean{LanglandsFirstMainLemma.PrimeCyclicExtension.lowerRamificationGroup_eq_top_of_le_break}
\lean{LanglandsFirstMainLemma.PrimeCyclicExtension.lowerRamificationGroup_eq_bot_of_break_lt}
\lean{LanglandsFirstMainLemma.PrimeCyclicExtension.lowerRamificationGroup_eq_bot_or_eq_top}
\lean{LanglandsFirstMainLemma.PrimeCyclicExtension.lowerRamificationGroup_eq_top_iff_le_break}
\lean{LanglandsFirstMainLemma.PrimeCyclicExtension.lowerRamificationGroup_eq_bot_iff_break_lt}
\lean{LanglandsFirstMainLemma.PrimeCyclicExtension.isLowerBreak_unique}
\lean{LanglandsFirstMainLemma.PrimeCyclicExtension.generator_mem_break_and_not_mem_succ}
\lean{LanglandsFirstMainLemma.PrimeCyclicExtension.existsUnique_isLowerBreak_of_zero_eq_top}
\lean{LanglandsFirstMainLemma.PrimeCyclicExtension.lowerBreakIndex}
\lean{LanglandsFirstMainLemma.PrimeCyclicExtension.lowerBreakIndex_isLowerBreak}
\lean{LanglandsFirstMainLemma.PrimeCyclicExtension.lowerBreakIndex_eq}
\leanok
\SourceText{Theorem thm:first-main and Section Ramification notation and the cyclic-prime norm theorem.}
\uses{mod:delta-firstmainstatement,mod:ramification-lowergroups}
\FormalizationContract The base structure's sole field reuses the existing
theorem-level cyclic-prime predicate, whose conjuncts are Galoisness,
cyclicity of the actual Galois group, and primality of the canonical
extension degree.  All choices and invariant identities are derived, and
the lower filtration is reused rather than copied.  A nonnegative break is
constructed only from the explicit hypothesis $G_0=G$; no ramifiedness,
total-ramification, tame--wild, Herbrand, different, trace, norm, or
conductor theorem is stored as structure data or proved in this node.
\end{definition}
